% p307_1; source PDF page 307; manually verified editable table.
\begin{tikzpicture}[x=1cm,y=-1cm,inner sep=0pt,outer sep=0pt]
\path[use as bounding box] (0,0) rectangle (16.000,9.760);
\draw[line width=.4pt] (0.000,0.000) rectangle (2.700,0.815);
\node[align=center,text width=2.440cm,font={\linespread{1}\fontsize{10}{13.50}\selectfont\bfseries}] at (1.350,0.407) {非法占有目的};
\draw[line width=.4pt] (2.700,0.000) rectangle (4.500,0.815);
\node[align=center,text width=1.540cm,font={\linespread{1}\fontsize{10}{13.50}\selectfont\bfseries}] at (3.600,0.407) {行为对象};
\draw[line width=.4pt] (4.500,0.000) rectangle (6.300,0.815);
\node[align=center,text width=1.540cm,font={\linespread{1}\fontsize{10}{13.50}\selectfont\bfseries}] at (5.400,0.407) {行为方式};
\draw[line width=.4pt] (6.300,0.000) rectangle (8.850,0.815);
\node[align=center,text width=2.290cm,font={\linespread{1}\fontsize{10}{13.50}\selectfont\bfseries}] at (7.575,0.407) {结论};
\draw[line width=.4pt] (8.850,0.000) rectangle (16.000,0.815);
\node[align=center,text width=6.890cm,font={\linespread{1}\fontsize{10}{13.50}\selectfont\bfseries}] at (12.425,0.407) {范例};
\draw[line width=.4pt] (0.000,0.815) rectangle (2.700,2.630);
\node[align=left,text width=2.440cm,font={\linespread{1}\fontsize{10}{13.50}\selectfont}] at (1.350,1.722) {一开始就不想付费};
\draw[line width=.4pt] (2.700,0.815) rectangle (4.500,2.630);
\node[align=center,text width=1.540cm,font={\linespread{1}\fontsize{10}{13.50}\selectfont}] at (3.600,1.722) {有形财物};
\draw[line width=.4pt] (4.500,0.815) rectangle (6.300,2.630);
\node[align=center,text width=1.540cm,font={\linespread{1}\fontsize{10}{13.50}\selectfont}] at (5.400,1.722) {欺骗手段};
\draw[line width=.4pt] (6.300,0.815) rectangle (8.850,2.630);
\node[align=center,text width=2.290cm,font={\linespread{1}\fontsize{10}{13.50}\selectfont}] at (7.575,1.722) {构成诈骗罪};
\draw[line width=.4pt] (8.850,0.815) rectangle (16.000,2.630);
\node[align=left,text width=6.890cm,font={\linespread{1}\fontsize{10}{13.50}\selectfont}] at (12.425,1.722) {甲一开始就不想付费，欺骗店家提供饭菜，吃完逃离。甲构成诈骗罪，行为对象是饭菜本身};
\draw[line width=.4pt] (0.000,2.630) rectangle (2.700,9.760);
\node[align=left,text width=2.440cm,font={\linespread{1}\fontsize{10}{13.50}\selectfont}] at (1.350,6.195) {一开始想付费，产生债务后，不想付费};
\draw[line width=.4pt] (2.700,2.630) rectangle (4.500,9.760);
\node[align=center,text width=1.540cm,font={\linespread{1}\fontsize{10}{13.50}\selectfont}] at (3.600,6.195) {财产性利益};
\draw[line width=.4pt] (4.500,2.630) rectangle (6.300,6.045);
\node[align=center,text width=1.540cm,font={\linespread{1}\fontsize{10}{13.50}\selectfont}] at (5.400,4.338) {欺骗手段};
\draw[line width=.4pt] (6.300,2.630) rectangle (8.850,6.045);
\node[align=center,text width=2.290cm,font={\linespread{1}\fontsize{10}{13.50}\selectfont}] at (7.575,4.338) {构成诈骗罪};
\draw[line width=.4pt] (8.850,2.630) rectangle (16.000,6.045);
\node[align=left,text width=6.890cm,font={\linespread{1}\fontsize{10}{13.50}\selectfont}] at (12.425,4.338) {例1，甲开始有付费意思，吃完后用假币支付餐费。甲构成诈骗罪。\\{}例2，甲开始有付费意思，吃完后不想付费，欺骗老板：我下次来一起付。老板相信，放走甲。老板放弃了当场实现债权。甲构成诈骗罪};
\draw[line width=.4pt] (4.500,6.045) rectangle (6.300,9.760);
\node[align=center,text width=1.540cm,font={\linespread{1}\fontsize{10}{13.50}\selectfont}] at (5.400,7.902) {平和手段};
\draw[line width=.4pt] (6.300,6.045) rectangle (8.850,9.760);
\node[align=left,text width=2.290cm,font={\linespread{1}\fontsize{10}{13.50}\selectfont}] at (7.575,7.902) {多数说认为不构成盗窃罪。少数说认为构成盗窃罪};
\draw[line width=.4pt] (8.850,6.045) rectangle (16.000,9.760);
\node[align=left,text width=6.890cm,font={\linespread{1}\fontsize{10}{13.50}\selectfont}] at (12.425,7.902) {例1，甲开始有付费意思，吃完后不想付钱，从后门偷偷溜走。存在观点展示。\\{}例2，甲开始有付费意思，吃完后不想付钱，欺骗老板：我送下朋友就回来结账，然后溜走。老板没有放弃当场实现债权。存在观点展示};
\end{tikzpicture}
